% Analytical example: entries are logical indices, never measured values.
\begin{figure}[t]
\centering\small
\begin{tabular}{c@{\quad}c}
Producer components & Consumer components\\
$P=\{0,1\}$ & $Q=\{2,3\}$\\[3pt]
\begin{tabular}{|rrrr|}\hline
0&1&2&3\\\hline
4&5&6&7\\\hline
8&9&10&11\\\hline
12&13&14&15\\\hline
\end{tabular}
&
\begin{tabular}{|rrrr|}\hline
0&4&8&12\\\hline
1&5&9&13\\\hline
2&6&10&14\\\hline
3&7&11&15\\\hline
\end{tabular}\\[5pt]
Logical index $4r+c$ & Write address $4c+r$
\end{tabular}
\par\smallskip
\begin{tabular}{lcc}\hline
Boundary & Packets & Producers/consumer\\\hline
$N=16$, $P\cup Q=\{0,1,2,3\}$ & 1 & 4\\
$N=64$, same $P,Q$ & 4 & 4\\\hline
\end{tabular}
\caption{Dependency-derived components and handoff. Each row denotes a
component, and each entry is a logical index. Every consumer requires one
value from each of four producer rows. Successor-aligned writeback makes its
row consecutive in storage. At $N=64$, the unchanged boundary relation
repeats independently for four settings of the high bits.}
\label{fig:hierarchy}
\end{figure}
